\section{Turing Machine --- Build --- Zone}
%%% generated by the blueprint pipeline from TCSlib.Complexity.TuringMachine.Build.Zone
\subsection*{Turing Machine --- Build --- Zone: Introduction}

Declarations of the Lean module \texttt{TCSlib.Complexity.TuringMachine.Build.Zone}.

\subsection*{Turing Machine --- Build --- Zone: Results}

\begin{definition}[Bits encoding an optional Boolean]
\lean{Turing.zoneCellBits}
\label{Turing.zoneCellBits}
\leanok
For an optional Boolean $v$, associate the pair whose first coordinate records whether
$v$ is present and whose second coordinate is its value when present, taking false when
it is absent.
\end{definition}

\begin{definition}[Optional zone cell]
\lean{Turing.zoneCellOf}
\label{Turing.zoneCellOf}
\leanok
For Boolean values $p$ and $d$, define a partial Boolean cell that is absent when $p$ is
false and equals $d$ when $p$ is true.
\end{definition}

\begin{theorem}[Reconstruction from cell bits]
\lean{Turing.zoneCellOf_bits}
\label{Turing.zoneCellOf_bits}
\leanok
\uses{Turing.zoneCellBits, Turing.zoneCellOf}
\difficulty{4}
For every optional Boolean cell value, encoding it by its presence bit and its Boolean
value (using false when absent), and then reconstructing the optional value from those
two bits, returns the original cell value.
\end{theorem}

\begin{definition}[Zone capacity]
\lean{Turing.zoneCapacity}
\label{Turing.zoneCapacity}
\leanok
For each natural number $i$, define the capacity to be $2^{i+1}$.
\end{definition}

\begin{definition}[Zone base]
\lean{Turing.zoneBase}
\label{Turing.zoneBase}
\leanok
For each natural number $i$, define the zone base to be twice the natural-number
difference between $2^i$ and $1$.
\end{definition}

\begin{theorem}[Successive zone bases]
\lean{Turing.zoneBase_succ}
\label{Turing.zoneBase_succ}
\leanok
\uses{Turing.zoneBase, Turing.zoneCapacity}
\difficulty{3}
For each natural number $i$, define the base of zone $i$ by $B_i=2(2^i-1)$ and its
capacity by $C_i=2\cdot 2^i$. Then $B_{i+1}=B_i+C_i$.
\end{theorem}

\begin{definition}[Zone index]
\lean{Turing.zoneIndex}
\label{Turing.zoneIndex}
\leanok
For each natural number $s$, define its zone index as the base-two logarithm, rounded
down, of $\log_2(s/2+1)$, where $s/2$ is integer division.
\end{definition}

\begin{theorem}[Characterization of the zone index]
\lean{Turing.zoneIndex_eq_iff}
\label{Turing.zoneIndex_eq_iff}
\leanok
\uses{Turing.zoneBase, Turing.zoneIndex}
\difficulty{3}
For all natural numbers $s$ and $i$, the zone index of $s$ equals $i$ if and only if
$2(2^i-1)\le s<2(2^{i+1}-1)$.
\end{theorem}

\begin{definition}[Zone contents]
\lean{Turing.ZoneContents}
\label{Turing.ZoneContents}
\leanok
\uses{Turing.zoneCapacity}
For a natural number $\ell$, a zone configuration consists of a possibly absent Boolean
home cell and two families of lists of possibly absent Boolean values, indexed by
$i<\ell$. Each list in either family has length at most $2^{i+1}$.
\end{definition}

\begin{definition}[Zone list slot]
\lean{Turing.zoneSlot}
\label{Turing.zoneSlot}
\leanok
\uses{Turing.zoneBase, Turing.zoneIndex}
Given a natural number $\ell$, a family of $\ell$ lists of possibly absent Boolean
values, and a natural number $s$, first set $i=\lfloor\log_2(\lfloor
s/2\rfloor+1)\rfloor$. If $i<\ell$, return the list entry at index $s-2(2^i-1)$ when
that index is within the list; otherwise return no value.
\end{definition}

\begin{definition}[Tape encoded by zone contents]
\lean{Turing.zoneTape}
\label{Turing.zoneTape}
\leanok
\uses{Turing.ZoneContents, Turing.zoneSlot}
Given zone contents with $\ell$ zones, define a tape on integer coordinates as follows.
At coordinates $0$ and $1$, return respectively whether the home cell is present and its
Boolean value, using false when it is absent. For $c\ge2$, let $n=c-2$; read the
right-zone slot selected by $\lfloor n/2\rfloor$, encoding its presence at even $n$ and
its Boolean value at odd $n$. For $c<0$, let $n=-c-1$; read the left-zone slot selected
by $\lfloor n/2\rfloor$, encoding its Boolean value at even $n$ and its presence at odd
$n$. In either zone, an unavailable slot gives no tape value.
\end{definition}

\begin{definition}[Empty zone contents]
\lean{Turing.ZoneContents.empty}
\label{Turing.ZoneContents.empty}
\leanok
\uses{Turing.ZoneContents}
For each natural number $\ell$, define the empty zone configuration to have an absent
home cell and empty lists in every left and right zone.
\end{definition}

\begin{theorem}[Tape contents of an empty zone]
\lean{Turing.zoneTape_empty}
\label{Turing.zoneTape_empty}
\leanok
\uses{Turing.ZoneContents.empty, Turing.zoneSlot, Turing.zoneTape}
\difficulty{1}
For any natural number $\ell$, the tape associated with empty zone contents has value
$\mathrm{some}(\mathrm{false})$ at positions $0$ and $1$, and is undefined at every
other integer position.
\end{theorem}

\begin{theorem}[Tape is blank outside the zones]
\lean{Turing.zoneTape_blank_outside}
\label{Turing.zoneTape_blank_outside}
\leanok
\uses{Turing.ZoneContents, Turing.zoneBase, Turing.zoneIndex, Turing.zoneIndex_eq_iff, Turing.zoneSlot, Turing.zoneTape}
\difficulty{6}
Let $z$ be zone contents with $\ell$ zones, and let $c$ be an integer satisfying
$2b_\ell+1<|c|$, where $b_\ell$ is the base of the zone following the last represented zone. Then the tape value at cell $c$ is blank.
the tape content at cell $c$ is blank.
\end{theorem}

\begin{definition}[Concatenated zone contents]
\lean{Turing.zoneSide}
\label{Turing.zoneSide}
\leanok
For a natural number $\ell$ and a family $w$ of lists indexed by $i=0,\ldots,\ell-1$,
form a single list by concatenating $w(0),w(1),\ldots,w(\ell-1)$ in that order. Each
list contains optional Boolean values.
\end{definition}

\begin{theorem}[Side word at successor length]
\lean{Turing.zoneSide_succ}
\label{Turing.zoneSide_succ}
\leanok
\uses{Turing.zoneSide}
\difficulty{1}
Let $n$ be a natural number and let $w$ assign a finite list of optional Boolean values
to each of $n+1$ indices. The concatenation of all these lists in index order is the
list at index $0$ followed by the concatenation of the lists at indices $1$ through $n$.
\end{theorem}

\begin{theorem}[Adjacent zone concatenation]
\lean{Turing.zoneSide_adjacent}
\label{Turing.zoneSide_adjacent}
\leanok
\uses{Turing.zoneSide, Turing.zoneSide_succ}
\difficulty{6}
Let $n$ be a natural number, let $a+1<n$, and let $w$ and $v$ assign a finite list of
optional Boolean values to each of the $n$ indices. If the concatenations of the lists
at indices $a$ and $a+1$ agree, and the lists assigned by $w$ and $v$ agree at every
other index, then the concatenations of all lists in index order agree.
\end{theorem}

\begin{definition}[Shift zone contents inward]
\lean{Turing.zoneShiftInW}
\label{Turing.zoneShiftInW}
\leanok
Let $\ell$ be a natural number, let $i$ be a natural number, and let $w$ assign a list
of optional Boolean values to each index from $0$ to $\ell-1$. If $1\le i<\ell$ and the
list at index $i-1$ is empty, split the list at index $i$ after its first $2^{i-1}$
entries, placing the initial segment at index $i-1$ and the remainder at index $i$. In
every other case, leave all lists unchanged.
\end{definition}

\begin{definition}[Shift zone contents outward]
\lean{Turing.zoneShiftOutW}
\label{Turing.zoneShiftOutW}
\leanok
\uses{Turing.zoneCapacity}
Let $\ell$ be a natural number, let $i$ be a natural number, and let $w$ assign a list
of optional Boolean values to each index from $0$ to $\ell-1$. If $1\le i<\ell$, the
list at index $i-1$ has length exactly $2\cdot 2^{i-1}$, and the list at index $i$ has
length at most $2\cdot 2^i-2^{i-1}$, split the list at index $i-1$ after its first
$2^{i-1}$ entries. Keep the initial segment at index $i$ and concatenate the remaining
segment before the existing list at index $i$. In every other case, leave all lists
unchanged.
\end{definition}

\begin{theorem}[Shift-in preserves concatenation]
\lean{Turing.zoneSide_shiftInW}
\label{Turing.zoneSide_shiftInW}
\leanok
\uses{Turing.zoneShiftInW, Turing.zoneSide, Turing.zoneSide_adjacent}
\difficulty{4}
Let $w$ assign a finite list of optional Boolean values to each of $\ell$ indices, and
let the side word be the concatenation of these lists in index order. For any natural
number $i$, shifting in leaves the side word unchanged: when $1\le i<\ell$ and the list
at index $i-1$ is empty, the list at index $i$ is split after its first $2^{i-1}$
entries and those two pieces are placed at indices $i-1$ and $i$; in every other case
the lists are unchanged.
\end{theorem}

\begin{theorem}[Shift-out preserves concatenation]
\lean{Turing.zoneSide_shiftOutW}
\label{Turing.zoneSide_shiftOutW}
\leanok
\uses{Turing.zoneCapacity, Turing.zoneShiftOutW, Turing.zoneSide, Turing.zoneSide_adjacent}
\difficulty{4}
Let $w$ assign a finite list of optional Boolean values to each of $\ell$ indices, and
let the side word be the concatenation of these lists in index order. For any natural
number $i$, shifting out leaves the side word unchanged: when $1\le i<\ell$, the list at
index $i-1$ has length exactly $C_{i-1}=2\cdot 2^{i-1}$, and the list at index $i$ has
room for $2^{i-1}$ more entries, the first $2^{i-1}$ entries of the list at $i-1$ are
placed there, followed by the remainder of that list and the original list at $i$; in
every other case the lists are unchanged.
\end{theorem}

\begin{theorem}[Inward shift with an empty donor]
\lean{Turing.zoneShiftInW_full_donor}
\label{Turing.zoneShiftInW_full_donor}
\leanok
\uses{Turing.zoneShiftInW}
\difficulty{3}
Let $\ell$ be a natural number, let $i$ satisfy $1\le i<\ell$, and assign a finite word
of optional Boolean values to each of the $\ell$ zones. If the word in zone $i-1$ is
empty, then after shifting inward at level $i$, the word in zone $i-1$ is the first
$2^{i-1}$ entries of the original word in zone $i$, and the word in zone $i$ is the
remainder after those entries are removed.
\end{theorem}

\begin{theorem}[Capacity under an inward shift]
\lean{Turing.zoneShiftInW_capacity}
\label{Turing.zoneShiftInW_capacity}
\leanok
\uses{Turing.zoneCapacity, Turing.zoneShiftInW}
\difficulty{4}
For every natural number $\ell$, every level $i\in\mathbb{N}$, and every assignment of
finite words of optional Boolean values to the $\ell$ zones, suppose each word has
length at most the capacity of its zone. After shifting the words inward at level $i$,
every word still has length at most the capacity of its zone.
\end{theorem}

\begin{definition}[Shift zone contents inward]
\lean{Turing.zoneShiftIn}
\label{Turing.zoneShiftIn}
\leanok
\uses{Turing.ZoneContents, Turing.zoneShiftInW, Turing.zoneShiftInW_capacity}
Let $\ell$ be a natural number and let $z$ consist of a home cell containing an optional
Boolean value, left and right lists of optional Boolean values indexed by
$0,\ldots,\ell-1$, with each list at index $i$ having length at most $2\cdot 2^i$. Given
a side and a natural number $i$, leave the home cell unchanged and shift the lists
inward at index $i$ on the side opposite the indicated side; leave the indicated side’s
lists unchanged.
\end{definition}

\begin{theorem}[Capacity under an outward shift]
\lean{Turing.zoneShiftOutW_capacity}
\label{Turing.zoneShiftOutW_capacity}
\leanok
\uses{Turing.zoneCapacity, Turing.zoneShiftOutW}
\difficulty{4}
For every natural number $\ell$, every level $i\in\mathbb{N}$, and every assignment of
finite words of optional Boolean values to the $\ell$ zones, suppose each word has
length at most the capacity of its zone. After shifting the words outward at level $i$,
every word still has length at most the capacity of its zone.
\end{theorem}

\begin{definition}[Shift zone contents outward]
\lean{Turing.zoneShiftOut}
\label{Turing.zoneShiftOut}
\leanok
\uses{Turing.ZoneContents, Turing.zoneShiftOutW, Turing.zoneShiftOutW_capacity}
Let $\ell$ be a natural number and let $z$ consist of a home cell containing an optional
Boolean value, left and right lists of optional Boolean values indexed by
$0,\ldots,\ell-1$, with each list at index $i$ having length at most $2\cdot 2^i$. Given
a side and a natural number $i$, leave the home cell unchanged and shift the lists
outward at index $i$ on the side opposite the indicated side; leave the indicated side’s
lists unchanged.
\end{definition}

\begin{definition}[Write the home cell]
\lean{Turing.zoneHomeWrite}
\label{Turing.zoneHomeWrite}
\leanok
\uses{Turing.ZoneContents}
Let $\ell$ be a natural number and let $z$ consist of a home cell containing an optional
Boolean value, left and right lists of optional Boolean values indexed by
$0,\ldots,\ell-1$, with each list at index $i$ having length at most $2\cdot 2^i$.
Replacing the home cell by any chosen optional Boolean value leaves both families of
lists unchanged.
\end{definition}

\begin{theorem}[Tape values away from the home cells]
\lean{Turing.zoneTape_homeWrite}
\label{Turing.zoneTape_homeWrite}
\leanok
\uses{Turing.ZoneContents, Turing.zoneHomeWrite, Turing.zoneTape}
\difficulty{1}
Let $\ell$ be a natural number, let $z$ be zone contents with $\ell$ zones, and let $v$
be an optional Boolean value. If $c$ is an integer with $c \ne 0$ and $c \ne 1$, then
changing the home value of $z$ to $v$ leaves the tape value at position $c$ unchanged.
\end{theorem}

\begin{definition}[Move the head right]
\lean{Turing.zoneMoveRight}
\label{Turing.zoneMoveRight}
\leanok
\uses{Turing.ZoneContents, Turing.zoneCapacity}
Let $\ell>0$, and let $z$ consist of a home cell containing an optional Boolean value
and left and right lists of optional Boolean values indexed by $0,\ldots,\ell-1$, each
of length at index $j$ at most $2\cdot 2^j$. Suppose the left list at index $0$ has room
for one more entry, so its length plus one is at most $2$. The new home cell is the head
of the right list at index $0$, interpreted as an optional Boolean value; prepend the
old home cell to the left list at index $0$ and remove the head of the right list at
index $0$. All other lists remain unchanged.
\end{definition}

\begin{definition}[Move the head left]
\lean{Turing.zoneMoveLeft}
\label{Turing.zoneMoveLeft}
\leanok
\uses{Turing.ZoneContents, Turing.zoneCapacity}
Let $\ell>0$, and let $z$ consist of a home cell containing an optional Boolean value
and left and right lists of optional Boolean values indexed by $0,\ldots,\ell-1$, each
of length at index $j$ at most $2\cdot 2^j$. Suppose the right list at index $0$ has
room for one more entry, so its length plus one is at most $2$. The new home cell is the
head of the left list at index $0$, interpreted as an optional Boolean value; remove the
head of the left list at index $0$ and prepend the old home cell to the right list at
index $0$. All other lists remain unchanged.
\end{definition}

\begin{definition}[Conditional zone move]
\lean{Turing.zoneMove}
\label{Turing.zoneMove}
\leanok
\uses{Turing.ZoneContents, Turing.zoneCapacity, Turing.zoneMoveLeft, Turing.zoneMoveRight}
For a zone configuration with $\ell$ zones, the zone-move operation leaves the
configuration unchanged if there are no zones. Otherwise, when moving right it performs
the rightward move if the first left-zone word can accommodate one more entry, and when
moving left it performs the leftward move if the first right-zone word can accommodate
one more entry; if the relevant word cannot accommodate that entry, it leaves the
configuration unchanged.
\end{definition}

\begin{theorem}[Rightward move updates side words]
\lean{Turing.zoneSide_moveRight}
\label{Turing.zoneSide_moveRight}
\leanok
\uses{Turing.ZoneContents, Turing.zoneCapacity, Turing.zoneMoveRight, Turing.zoneSide, Turing.zoneSide_succ}
\difficulty{4}
Let $\ell>0$, set $C_i=2\cdot 2^i$, and let $z$ consist of a home value and left and
right lists at each index $i<\ell$, each of length at most $C_i$. Suppose the left list
at index $0$ has length plus one at most $C_0$, and the right list at index $0$ is
nonempty. After moving the configuration one position to the right, the concatenation of
the left lists is the home value prepended to its original concatenation, while the
concatenation of the right lists is the tail of its original concatenation.
\end{theorem}

\begin{definition}[Zone shift pair]
\lean{Turing.zoneStepPair}
\label{Turing.zoneStepPair}
\leanok
\uses{Turing.ZoneContents, Turing.zoneShiftIn, Turing.zoneShiftOut}
For a zone configuration and a natural-number index $i$, the zone-step operation first
shifts into zone $i$ in the rightward direction and then shifts out of zone $i$ in the
leftward direction.
\end{definition}

\begin{definition}[Rightward zone cascade]
\lean{Turing.zoneCascadeRight}
\label{Turing.zoneCascadeRight}
\leanok
\uses{Turing.ZoneContents, Turing.zoneMove, Turing.zoneStepPair}
For a zone configuration and a natural number $j$, the rightward cascade applies the
zone-step operation at indices $j,j-1,\ldots,1$ in that order, performs a rightward zone
move, and then applies the zone-step operation at indices $1,2,\ldots,j$ in that order.
When $j=0$, both sequences of zone steps are empty.
\end{definition}

\begin{theorem}[Paired shift preserves side words]
\lean{Turing.zoneStepPair_side}
\label{Turing.zoneStepPair_side}
\leanok
\uses{Turing.ZoneContents, Turing.zoneSide, Turing.zoneSide_shiftInW, Turing.zoneSide_shiftOutW, Turing.zoneStepPair}
\difficulty{1}
For any zone contents $z$ with $\ell$ zones and any natural number $i$, the concatenation
of all left zone words is unchanged by the paired shift at zone $i$, and the
concatenation of all right zone words is unchanged as well.
\end{theorem}

\begin{theorem}[Paired shift frame]
\lean{Turing.zoneStepPair_frame}
\label{Turing.zoneStepPair_frame}
\leanok
\uses{Turing.ZoneContents, Turing.zoneShiftInW, Turing.zoneShiftOutW, Turing.zoneStepPair}
\difficulty{3}
Let $z$ be zone contents with $\ell$ zones, and let $c$ be an integer satisfying
$2b_\ell+1<|c|$, where $b_\ell$ is the base of the zone following the last represented zone. Then the tape value at cell $c$ is blank.
left and right words at zone $k$ are unchanged.
\end{theorem}

\begin{theorem}[Active paired zone shift]
\lean{Turing.zoneStepPair_active}
\label{Turing.zoneStepPair_active}
\leanok
\uses{Turing.ZoneContents, Turing.zoneCapacity, Turing.zoneShiftIn, Turing.zoneShiftInW, Turing.zoneShiftOut, Turing.zoneShiftOutW, Turing.zoneStepPair}
\difficulty{1}
Let $\ell$ be a natural number, let $i+1<\ell$, and let $z$ be zone contents for $\ell$
zones. Suppose the right word in zone $i$ is empty, the left word in zone $i$ has length
equal to its capacity, and the left word in zone $i+1$ has length plus $2^i$ at most its
capacity. After the paired shift at zone $i+1$, the right word in zone $i$ is the first
$2^i$ symbols of the original right word in zone $i+1$, the left word in zone $i$ is the
first $2^i$ symbols of the original left word in zone $i$, the right word in zone $i+1$
is the remainder of the original right word there, and the left word in zone $i+1$ is
the remainder of the original left word in zone $i$ followed by the original left word
in zone $i+1$.
\end{theorem}

\begin{theorem}[Successive rightward cascades]
\lean{Turing.zoneCascadeRight_succ}
\label{Turing.zoneCascadeRight_succ}
\leanok
\uses{Turing.ZoneContents, Turing.zoneCascadeRight, Turing.zoneStepPair}
\difficulty{1}
For every natural number $\ell$, every $j\in\mathbb{N}$, and every configuration of
$\ell$ zones, the rightward cascade of length $j+1$ is obtained by applying one paired
zone shift at level $j+1$, then the length-$j$ rightward cascade, then another paired
shift at level $j+1$.
\end{theorem}

\begin{theorem}[Cascade leaves higher zones unchanged]
\lean{Turing.zoneCascadeRight_frame}
\label{Turing.zoneCascadeRight_frame}
\leanok
\uses{Turing.ZoneContents, Turing.zoneCascadeRight, Turing.zoneCascadeRight_succ, Turing.zoneMove, Turing.zoneMoveRight, Turing.zoneStepPair_frame}
\difficulty{4}
Let $\ell$ be a natural number, let $j$ be a natural number, and let $z$ be zone
contents with $\ell$ levels. For any level $k<\ell$ with $j<k$, the rightward cascade
through $j$ levels leaves both the left and right words at level $k$ unchanged.
\end{theorem}

\begin{theorem}[Rightward cascade preserves side words]
\lean{Turing.zoneSide_cascadeRight}
\label{Turing.zoneSide_cascadeRight}
\leanok
\uses{Turing.ZoneContents, Turing.zoneCapacity, Turing.zoneCascadeRight, Turing.zoneCascadeRight_succ, Turing.zoneMove, Turing.zoneSide, Turing.zoneSide_moveRight, Turing.zoneStepPair, Turing.zoneStepPair_active, Turing.zoneStepPair_frame, Turing.zoneStepPair_side}
\difficulty{5}
For a natural number $\ell$, let $C_i=2\cdot 2^i$ be the capacity of zone $i$, and let
$z$ be a configuration of a home value and left and right lists at each index $i<\ell$,
with each list having length at most $C_i$. Let $j<\ell$. Suppose the right lists at all
indices $k<j$ are empty, the left list at each such index has length $C_k$, the right
list at index $j$ is nonempty, and the left list at index $j$ has length at most
$C_j-2^{j-1}$. After the rightward cascade through zone $j$, the concatenation of the
left lists is the home value prepended to its original concatenation, and the
concatenation of the right lists is the tail of its original concatenation.
\end{theorem}

\begin{theorem}[Zone lengths after a rightward cascade]
\lean{Turing.zoneCascadeRight_lengths}
\label{Turing.zoneCascadeRight_lengths}
\leanok
\uses{Turing.ZoneContents, Turing.zoneCapacity, Turing.zoneCascadeRight, Turing.zoneCascadeRight_frame, Turing.zoneCascadeRight_succ, Turing.zoneMove, Turing.zoneMoveRight, Turing.zoneStepPair, Turing.zoneStepPair_active, Turing.zoneStepPair_frame}
\difficulty{6}
Let $\ell$ be a natural number, let $j<\ell$, and let $z$ be zone contents with $\ell$
levels. Suppose the right words at all levels $k<j$ are empty, the left words at those
levels fill their zones to capacity, the right word at level $j$ has length at least
$2^j$, and the left word at level $j$ has room for at least $2^j$ additional symbols.
After the rightward cascade through $j$ levels, each level $k<j$ has left and right word
lengths both equal to $2^k$; at level $j$, the right word's length is its original
length minus $2^j$, and the left word's length is its original length plus $2^j$.
\end{theorem}

\begin{theorem}[Zero-step rightward cascade]
\lean{Turing.zoneCascadeRight_zero}
\label{Turing.zoneCascadeRight_zero}
\leanok
\uses{Turing.ZoneContents, Turing.zoneCapacity, Turing.zoneCascadeRight, Turing.zoneMove, Turing.zoneSide, Turing.zoneSide_moveRight}
\difficulty{1}
Let $\ell>0$ and let $z$ be a configuration of $\ell$ zones. If the innermost right zone
is nonempty and the innermost left zone has room for one additional cell, then the
concatenated left-side contents after the zero-step rightward cascade equal the home
cell followed by the original concatenated left-side contents, and the concatenated
right-side contents equal the tail of the original concatenated right-side contents.
\end{theorem}

\begin{theorem}[Blocked rightward cascade preserves contents]
\lean{Turing.zoneCascadeRight_blocked}
\label{Turing.zoneCascadeRight_blocked}
\leanok
\uses{Turing.ZoneContents, Turing.zoneCapacity, Turing.zoneCascadeRight, Turing.zoneCascadeRight_succ, Turing.zoneMove, Turing.zoneShiftIn, Turing.zoneShiftOut, Turing.zoneShiftOutW, Turing.zoneSide, Turing.zoneStepPair, Turing.zoneStepPair_side}
\difficulty{5}
Let $\ell$ be a natural number, let $j<\ell$, and consider zone contents with left and
right words at each of the $\ell$ levels. Suppose each left word at a level $k<j$ fills
that zone to capacity, and the left word at level $j$ also fills its zone to capacity.
Then the rightward cascade through $j$ levels preserves the concatenated left-side word
and the concatenated right-side word.
\end{theorem}

\begin{theorem}[Rightward cascade cost bound]
\lean{Turing.zoneCascade_cost_le}
\label{Turing.zoneCascade_cost_le}
\leanok
\difficulty{4}
For every $j\in\mathbb{N}$, the sum over $0\le i<j$ of $4\bigl(2^{i+1}+i+2\bigr)$ is at
most $16\cdot 2^j$.
\end{theorem}

\begin{definition}[Encoding a zone word]
\lean{Turing.zoneStageWord}
\label{Turing.zoneStageWord}
\leanok
Given a side and a word whose entries are either absent or Boolean values, encode each
entry by two bits. For the right side, the pair records first whether the entry is
present and then its value, taking an absent value to be false; for the other side,
record the value first and presence second. Concatenate these pairs in the original
order.
\end{definition}

\begin{theorem}[Zone stage word length]
\lean{Turing.zoneStageWord_length}
\label{Turing.zoneStageWord_length}
\leanok
\uses{Turing.zoneStageWord}
\difficulty{4}
For either choice of side and every finite word $w$ of optional Boolean symbols, the
associated stage word has length $2|w|$.
\end{theorem}

\begin{theorem}[Entries of the encoded zone word]
\lean{Turing.zoneStageWord_getElem}
\label{Turing.zoneStageWord_getElem}
\leanok
\uses{Turing.zoneStageWord}
\difficulty{5}
Let $w$ be a finite list of optional Boolean values, and let $b$ be a Boolean choosing
an encoding order. Each list entry is encoded by two Boolean entries: when $b$ is true,
they are its presence indicator followed by its value, and when $b$ is false, they are
its value followed by its presence indicator. An absent value is encoded using false for
its value. For every position $p$, the encoded word has an entry there exactly when $w$
has an entry at position $\lfloor p/2\rfloor$; the encoded entry is the first of that
pair when $p$ is even and the second when $p$ is odd.
\end{theorem}

\begin{theorem}[Slot within a zone]
\lean{Turing.zoneStageSlot}
\label{Turing.zoneStageSlot}
\leanok
\uses{Turing.zoneBase, Turing.zoneBase_succ, Turing.zoneCapacity, Turing.zoneIndex, Turing.zoneIndex_eq_iff, Turing.zoneSlot}
\difficulty{4}
For each natural number $i$, set $B_i=2(2^i-1)$ and assign index $s$ the zone index
$I(s)=\lfloor\log_2(s/2+1)\rfloor$, with integer division understood. Let $w$ assign a
finite list of optional Boolean values to each of $\ell$ indices. For any index $i<\ell$
and any natural number $p<2\cdot 2^i$, the entry at position $B_i+p$ in the slot
assignment is exactly the entry at position $p$ of the list assigned by $w$ to index
$i$, with an out-of-range list position interpreted as absent.
\end{theorem}

\begin{theorem}[Right zone window]
\lean{Turing.zoneStage_rightWindow}
\label{Turing.zoneStage_rightWindow}
\leanok
\uses{Turing.FinTM.bufferTape, Turing.FinTM.bufferTape_nat, Turing.ZoneContents, Turing.zoneBase, Turing.zoneCapacity, Turing.zoneStageSlot, Turing.zoneStageWord, Turing.zoneStageWord_getElem, Turing.zoneTape}
\difficulty{5}
Let $z$ be zone contents with $\ell$ zones, and let $c$ be an integer satisfying
$2b_\ell+1<|c|$, where $b_\ell$ is the base of the zone following the last represented zone. Then the tape value at cell $c$ is blank.
the right word at zone $i$ by the pair consisting of whether it is present followed by
its default Boolean value. At the integer cell $2b_i+2+p$, where $b_i$ is the base of
zone $i$, the tape content equals the content at position $p$ of this encoded word
stored from position zero, with blank content beyond the word.
\end{theorem}

\begin{theorem}[Left zone window]
\lean{Turing.zoneStage_leftWindow}
\label{Turing.zoneStage_leftWindow}
\leanok
\uses{Turing.FinTM.bufferTape, Turing.FinTM.bufferTape_nat, Turing.ZoneContents, Turing.zoneBase, Turing.zoneCapacity, Turing.zoneStageSlot, Turing.zoneStageWord, Turing.zoneStageWord_getElem, Turing.zoneTape}
\difficulty{5}
Let $z$ be zone contents with $\ell$ zones, and let $c$ be an integer satisfying
$2b_\ell+1<|c|$, where $b_\ell$ is the base of the zone following the last represented zone. Then the tape value at cell $c$ is blank.
the left word at zone $i$ by the pair consisting of its default Boolean value followed
by whether it is present. At the integer cell $-(2b_i)-1-p$, where $b_i$ is the base of
zone $i$, the tape content equals the content at position $p$ of this encoded word
stored from position zero, with blank content beyond the word.
\end{theorem}

\begin{theorem}[Zone window bounds]
\lean{Turing.zoneStage_window_bounds}
\label{Turing.zoneStage_window_bounds}
\leanok
\uses{Turing.zoneBase, Turing.zoneBase_succ, Turing.zoneCapacity}
\difficulty{3}
For a zone index $i$ and an integer $p$ with $-1\le p\le 2$ times the capacity of zone
$i$, both $2b_i+2+p$ and $-(2b_i)-1-p$ lie in the integer interval from $-(2b_{i+1}+2)$
to $2b_{i+1}+2$, where $b_j$ denotes the base of zone $j$.
\end{theorem}

\begin{theorem}[Bounded implementation of an inward zone shift]
\lean{Turing.FinTM.exists_zoneShiftInTM}
\label{Turing.FinTM.exists_zoneShiftInTM}
\leanok
\uses{Turing.Cfg, Turing.FinTM, Turing.FinTM.bufferTape, Turing.MultiTapeTM.runFrom, Turing.MultiTapeTM.spaceUsedByTape, Turing.ZoneContents, Turing.zoneBase, Turing.zoneShiftIn, Turing.zoneTape}
For each direction $s\in\{\mathrm{false},\mathrm{true}\}$, there are a finite-state
Turing machine $Z$ over Boolean symbols and a constant $c\in\mathbb N$ such that $Z$ has
two work tapes and the following holds. For all natural numbers $\ell,i$ with $1\le
i<\ell$, every zone content $z$ with $\ell$ zones, every Boolean input word $x$, and
every initial configuration $d$ whose state is the initial state of $Z$, whose first
work tape stores the encoding of $z$, whose second stores $i$ copies of $\mathrm{true}$,
whose work-tape heads are both at zero, and whose output is empty, there is a time $T\le
c(2^i+i+1)$ at which the machine has halted, its first tape stores the encoding of the
inward-shifted zone content, its second tape still stores $i$ copies of $\mathrm{true}$,
both work-tape heads are at zero, its output is empty, and its input-head position is
unchanged. At every time $t<T$ it is not halted; throughout times $t\le T$, the first
work-tape head lies in $[-(2B+2),2B+2]$, where $B$ is the base of zone $i+1$; and the
second tape visits at most $c(2^i+i+1)$ cells by time $T$.
\end{theorem}

\begin{theorem}[Bounded implementation of an outward zone shift]
\lean{Turing.FinTM.exists_zoneShiftOutTM}
\label{Turing.FinTM.exists_zoneShiftOutTM}
\leanok
\uses{Turing.Cfg, Turing.FinTM, Turing.FinTM.bufferTape, Turing.MultiTapeTM.runFrom, Turing.MultiTapeTM.spaceUsedByTape, Turing.ZoneContents, Turing.zoneBase, Turing.zoneShiftOut, Turing.zoneTape}
For each direction $s\in\{\mathrm{false},\mathrm{true}\}$, there are a finite-state
Turing machine $Z$ over Boolean symbols and a constant $c\in\mathbb N$ such that $Z$ has
two work tapes and the following holds. For all natural numbers $\ell,i$ with $1\le
i<\ell$, every zone content $z$ with $\ell$ zones, every Boolean input word $x$, and
every initial configuration $d$ whose state is the initial state of $Z$, whose first
work tape stores the encoding of $z$, whose second stores $i$ copies of $\mathrm{true}$,
whose work-tape heads are both at zero, and whose output is empty, there is a time $T\le
c(2^i+i+1)$ at which the machine has halted, its first tape stores the encoding of the
outward-shifted zone content, its second tape still stores $i$ copies of
$\mathrm{true}$, both work-tape heads are at zero, its output is empty, and its
input-head position is unchanged. At every time $t<T$ it is not halted; throughout times
$t\le T$, the first work-tape head lies in $[-(2B+2),2B+2]$, where $B$ is the base of
zone $i+1$; and the second tape visits at most $c(2^i+i+1)$ cells by time $T$.
\end{theorem}
